\markright{\mititulo \hspace{5cm} \thepage}
\pagestyle{fancy}

\section*{Resumen}
\addcontentsline{toc}{section}{Resumen}

\noindent Se estudia la variación anual de las horas de salida y puesta del Sol, y la
duración del día, en cinco ciudades ubicadas entre las latitudes $+60{,}17^\circ$\,N
(Helsinki, Finlandia) y $-53{,}16^\circ$\,S (Punta Arenas, Chile), para 19 fechas
distribuidas a lo largo del año 2026 en intervalos de veinte días.
Las efemérides solares se calcularon con el paquete \texttt{astroplan}
\citep{Morris2018}, que implementa el cálculo riguroso de eventos astronómicos mediante
la corrección estándar de horizonte ($-0{,}833^\circ$) debida a la refracción
atmosférica y al semidiámetro solar \citep{Meeus1998}.
Los resultados se validaron con el servicio \textit{JPL Horizons} de la NASA
\citep{Giorgini1996}, confirmando la coherencia de la metodología.
La duración del día en Helsinki varía desde $5$\,h\,$52$\,min (diciembre) hasta
$18$\,h\,$49$\,min (junio), mientras que en la Universidad de Antioquia
($\phi = +6{,}27^\circ$\,N) la variación es de apenas $\approx 43$\,min a lo largo del
año.
Los resultados confirman que la inclinación axial de la Tierra ($\varepsilon \approx
23{,}44^\circ$) es la causa directa de las estaciones del año y de las diferencias
asimétricas de insolación entre los hemisferios.

\vspace{1cm}\textit{Palabras clave:} eclíptica, declinación solar, horas de luz,
estaciones del año, hemisferios, latitud, \texttt{astroplan}, JPL Horizons

\newpage

\section*{Abstract}
\addcontentsline{toc}{section}{Abstract}

\noindent The annual variation of sunrise and sunset times, and day length, is studied
for five cities located between latitudes $+60.17^\circ$\,N (Helsinki, Finland) and
$-53.16^\circ$\,S (Punta Arenas, Chile), for 19 dates distributed throughout 2026 at
twenty-day intervals.
Solar ephemerides were computed with the \texttt{astroplan} package
\citep{Morris2018}, which implements rigorous calculation of astronomical events using
the standard horizon correction ($-0.833^\circ$) for atmospheric refraction and solar
semi-diameter \citep{Meeus1998}.
Results were validated against NASA's \textit{JPL Horizons} service
\citep{Giorgini1996}, confirming the consistency of the methodology.
Day length in Helsinki ranges from $5$\,h\,$52$\,min (December) to $18$\,h\,$49$\,min
(June), while at the Universidad de Antioquia ($\phi = +6.27^\circ$\,N) the variation
is only $\approx 43$\,min throughout the year.
The results confirm that the Earth's axial tilt ($\varepsilon \approx 23.44^\circ$) is
the direct cause of the seasons and of the asymmetric insolation differences between
hemispheres.

\vspace{1cm}\textit{Keywords:} ecliptic, solar declination, daylight hours, seasons,
hemispheres, latitude, \texttt{astroplan}, JPL Horizons

\newpage
